\chapter{Relative arithmetic cohomology}

\begin{lemma}[Extensions of flat modules]
\label{lem:relative-flat-extension}
\lean{EnvelopeProperty.Vanishing.RelativeArithmeticCohomologyFlat.flat_of_flat_submodule}
\leanok
Let $N$ be a submodule of an $A$-module $M$. If $N$ and $M/N$ are flat over $A$,
then $M$ is flat over $A$.
\end{lemma}
\begin{proof}
For an injection $P\to Q$, tensor the short exact sequence
$0\to N\to M\to M/N\to0$ with $P$ and $Q$. Both resulting rows are exact,
including at their left ends, because $M/N$ is flat. An element of
$P\otimes M$ killed in $Q\otimes M$ maps to zero in $P\otimes(M/N)$,
because $P\otimes(M/N)\to Q\otimes(M/N)$ is injective. It therefore lifts
to $P\otimes N$. Its image in $Q\otimes N$ is zero by left exactness of
the lower row. Flatness of $N$ makes its original lift zero. Thus
$P\otimes M\to Q\otimes M$ is injective for every injection, proving flatness.
\end{proof}

\begin{lemma}[Flat opcycles]
\label{lem:relative-flat-opcycles}
\uses{lem:relative-flat-extension}
\lean{EnvelopeProperty.Vanishing.RelativeArithmeticCohomologyFlat.opcycles_flat}
\leanok
Consider a complex $P\xrightarrow f M\xrightarrow g Q$. If its homology,
$Q$, and $\operatorname{coker}g$ are flat over $A$, then
$\operatorname{coker}f$ is flat over $A$.
\end{lemma}
\begin{proof}
The image of $g$ is flat, being a submodule of the flat module $Q$ with
flat quotient. The induced map $\operatorname{coker}f\to Q$ has kernel
equal to the original homology and image equal to $\operatorname{im}g$.
Apply the preceding lemma to this kernel and the resulting quotient.
\end{proof}

\begin{theorem}[Homology base change without finite arithmetic rank]
\label{thm:relative-flat-homology-base-change}
\uses{lem:relative-flat-opcycles}
\lean{EnvelopeProperty.Vanishing.RelativeArithmeticCohomologyBaseChange.cokernel_flat,EnvelopeProperty.Vanishing.RelativeArithmeticCohomologyBaseChange.preserves_kernel_of_flat_cokernel,EnvelopeProperty.Vanishing.RelativeArithmeticCohomologyBaseChange.shortHomologyIso,EnvelopeProperty.Vanishing.RelativeArithmeticCohomologyBaseChange.naturalHomologyIso}
\leanok
Let $K^\bullet$ be a nonnegative cochain complex of flat $A$-modules.
Suppose every $H^i(K^\bullet)$ is flat over $A$, and fix $N$ with
$K^{N+1}=0$. Then for every $A$-algebra $C$ and $0<i\leq N$ the canonical map
$C\otimes_A H^i(K^\bullet)\to H^i(C\otimes_A K^\bullet)$ is an
isomorphism.
\end{theorem}
\begin{proof}
The cokernel of $d^N$ is zero. Descending induction using the preceding
lemma shows that the cokernel of $d^i$ is flat for every $i\leq N$.
The kernel of each such differential therefore commutes with arbitrary tensor products:
factor it through its image, and tensor the two short exact sequences,
whose quotients are flat. Tensor products always preserve cokernels.
The description of homology as the cokernel of the map into cycles now
gives the asserted canonical isomorphism.
\end{proof}

\begin{lemma}[The same two coefficient actions]
\label{lem:relative-arithmetic-cech-coefficients}
\lean{EnvelopeProperty.Vanishing.RelativeArithmeticCohomologyCoefficients.openSectionsIso,EnvelopeProperty.Vanishing.RelativeArithmeticCohomologyCoefficients.termEquiv,EnvelopeProperty.Vanishing.RelativeArithmeticCohomologyCoefficients.termEquiv_coordinate,EnvelopeProperty.Vanishing.RelativeArithmeticCohomologyCoefficients.termEquiv_differential,EnvelopeProperty.Vanishing.RelativeArithmeticCohomologyCoefficients.complexIso,EnvelopeProperty.Vanishing.RelativeArithmeticCohomologyCoefficients.homologyIso}
\leanok
Let $A\to B$ be a homomorphism, $g:Y\to\operatorname{Spec}B$ a morphism,
and $M$ a module on $Y$. On the same ordered cover, the Čech complex with
its original $A$-action is canonically isomorphic to restriction of the
Čech complex with its original $B$-action. The isomorphism fixes every
open-section coordinate and commutes with every differential.
\end{lemma}
\begin{proof}
For each open $U$, both section modules have carrier $\Gamma(U,M)$.
The two $A$-actions agree because the structural map to
$\operatorname{Spec}A$ is the composite of $g$ and
$\operatorname{Spec}B\to\operatorname{Spec}A$; naturality of the
global-sections isomorphism for affine spectra gives the same coefficient
map $A\to\Gamma(U,\mathcal O_Y)$. The identity maps therefore give linear
isomorphisms. They commute with restrictions because both restriction
maps are those of the same original $M$. Taking products over the ordered
intersections gives the degreewise comparisons, and their compatibility
with the alternating sums of restrictions gives the complex comparison.
Restriction of scalars is exact, giving the asserted homology comparison.
\end{proof}

\begin{lemma}[Relative finite homology and arithmetic flat terms]
\label{lem:relative-arithmetic-coordinate-complex}
\uses{lem:relative-arithmetic-cech-coefficients}
\lean{EnvelopeProperty.Vanishing.RelativeArithmeticCohomologyCech.exists_coordinate_complex}
\leanok
Suppose $B$ is noetherian, $g:Y\to\operatorname{Spec}B$ is projective,
the composite $Y\to\operatorname{Spec}A$ is flat, and $M$ is invertible.
There is a finite ordered affine cover of $Y$ such that its Čech complex
has flat terms over $A$, finite homology over $B$, and terms zero above
the cardinality of the cover.
\end{lemma}
\begin{proof}
Choose a closed projective embedding over $B$ and restrict the standard
coordinate affine cover. Its intersections are affine, since projectivity
over the affine base makes $Y$ separated. The coordinate Čech finiteness
theorem applies to the coherent invertible module and yields homology
finite over $B$. Flatness of $Y/A$ and invertibility of $M$ make its
sections on each affine intersection flat over $A$. Finite products
preserve flatness, so the terms are flat. A strictly increasing tuple
of length exceeding the number of cover members does not exist, making
every such term zero. The preceding lemma identifies these two coefficient
structures on the same complex.
\end{proof}

\begin{lemma}[The original scalar tower comparison]
\label{lem:relative-arithmetic-scalar-tower-complex}
\lean{EnvelopeProperty.Vanishing.RelativeArithmeticCohomologyLocalization.iteratedShortIso}
\leanok
For compatible homomorphisms $A\to T\to C$, iterated scalar extension of
an original three-term complex is canonically isomorphic to its direct
scalar extension from $A$ to $C$.
\end{lemma}
\begin{proof}
On every original module $V$, the canonical map
$C\otimes_T(T\otimes_A V)\to C\otimes_A V$ sends
$c\otimes(t\otimes v)$ to $ct\otimes v$. Its inverse sends
$c\otimes v$ to $c\otimes(1\otimes v)$. Tensor balancing verifies that
both maps are well defined; their composites are identities on pure
tensors and hence on all tensors by additivity. Both maps commute with
each original differential, as follows by applying the differential to
$v$ in these same pure-tensor formulas. They therefore give the stated
whole-complex isomorphism.
\end{proof}

\begin{lemma}[Flat scalar extension of native homology]
\label{lem:relative-arithmetic-flat-scalar-homology}
\uses{thm:relative-flat-homology-base-change}
\lean{EnvelopeProperty.Vanishing.RelativeArithmeticCohomologyLocalization.flatHomologyIso}
\leanok
If $T$ is flat over $A$, then native homology of any original nonnegative
cochain complex $K$ commutes with scalar extension to $T$, with its
literal original coefficient action.
\end{lemma}
\begin{proof}
Tensoring with the flat $A$-module $T$ preserves kernels. Tensor right
exactness preserves cokernels. Apply these two facts to the original
three-term complex in each degree, describing homology as the cokernel
of the original incoming map into its kernel of the outgoing map.
Every map is the scalar extension of the same original linear map,
so the resulting comparison is the native homology comparison.
\end{proof}

\begin{theorem}[A free scalar fibre detects native homology]
\label{thm:relative-arithmetic-free-scalar-homology}
\uses{thm:relative-flat-homology-base-change,lem:relative-arithmetic-scalar-tower-complex,lem:kn-arithmetic-free-fibre-detection}
\lean{EnvelopeProperty.Vanishing.RelativeArithmeticCohomologyLocalization.natural_homology_isZero_of_free_scalar_fibre}
\leanok
Let $K$ be a nonnegative complex of flat $A$-modules, with $K^{N+1}=0$.
Suppose that all native homology of $T\otimes_A K$ is free over $T$.
For compatible scalar towers $A\to T\to C$ and $A\to T\to F$ with
$C\ne0$, if $H^j(C\otimes_A K)=0$ for $0<j\leq N$, then
$H^j(F\otimes_A K)=0$.
\end{theorem}
\begin{proof}
The extended complex has flat terms over $T$, flat homology because its
homology is free, and zero term in degree $N+1$. Bounded homology base
change and the scalar tower comparison identify the stipulated zero
group with $C\otimes_T H^j(T\otimes_A K)$. A free module whose extension
to the nonzero ring $C$ is zero has empty basis and is zero. Its extension
to $F$ is therefore zero. Apply the same bounded homology and scalar
tower comparisons to identify this zero module with the asserted actual
$F$-scalar homology.
\end{proof}

\begin{lemma}[The actual Čech free Witt point]
\label{lem:relative-arithmetic-cech-free-witt-point}
\uses{lem:relative-arithmetic-cech-coefficients,lem:relative-arithmetic-flat-scalar-homology,lem:kn-relative-arithmetic-simultaneous-free-witt-point}
\lean{EnvelopeProperty.Vanishing.RelativeArithmeticCohomologyFreeStage.exists_free_witt_point}
\leanok
Let $A$ be a characteristic-zero noetherian domain smooth and of finite
type over $\mathbf Z$, and let $B$ be a finite-type $A$-algebra. For an
actual bounded ordered Čech complex with homology finite over $B$, and
any integer $b$, there is a closed point $P$, a prime $\ell>b$, a finite
quotient field $A/P$ of characteristic $\ell$, and an actual lift
$A\to W_2(A/P)$ of the quotient map, such that all native homology
modules of the same complex extended to $A_P$ are free over $A_P$.
\end{lemma}
\begin{proof}
Restrict the relative $B$-module homology to $A$. Apply simultaneous
generic freeness and the large-characteristic Witt-point construction
to the finite family of homology modules in degrees at most the complex
bound. Localization is flat, so native homology of the same complex
extended to $A_P$ is its original homology tensored with $A_P$.
The original two-coefficient comparison and the universal tensor
description of localized modules identify this with the produced free
localization of the corresponding $B$-module homology. Beyond the bound,
the middle Čech term is zero, so its localized homology is zero and free.
Every scalar map and the Witt reduction equation come from the actual
simultaneous free-point construction.
\end{proof}
